\section{Scalability and Limitations}
\label{sec:discussion-limitations}

A primary limitation of neural Lyapunov certification lies in the trade-off between the network's expressive power,
the size of the architecture, and the dimensionality of the physical state-space.
While linear bound propagation (LiRPA) in $\alpha,\beta$-CROWN runs efficiently on GPUs, 
increasing the depth and width of neural networks significantly degrades verifiability. 
Layer-wise bound propagation accumulates relaxation errors across successive layers, 
while larger models introduce more unstable activation states that can substantially expand 
the size of the branch-and-bound search tree. 
Consequently, existing neural Lyapunov approaches rely on compact network architectures 
in order to maintain tight bounds~\citeTab{3}{yangLyapunovstableNeuralControl2024}.
Simultaneously, the computational complexity is constrained by the continuous state dimension $n_x$. 
Since formal certification over the region of interest $\mathcal{X}_C$ requires recursive input domain partitioning, 
this leads, in the worst case, to an exponential dependence on the state dimension.
The combination of the accumulation of neural relaxation errors and the curse of dimensionality 
limits rigorous formal verification to systems with low to medium state dimensionality
and compact neural representations. 

Furthermore, the pipeline fundamentally relies on the coverage of the expert imitation dataset $\mathcal{D}_\pi$. 
For the cartpole, approximately 30\,\% of the sampled trajectories are feasible. 
Consequently, the solver fails to identify feasible solutions for certain regions of the state space, 
such that these regions are only weakly represented by expert supervision.

Finally, the verification guarantees rely on a nominal discrete-time model $x_{k+1} = f(x_k, u_k)$ discretized with 
explicit Euler. While higher-order methods such as 4th-order Runge-Kutta (RK4) provide higher simulation accuracy, 
incorporating four intermediate forward passes into the computational graph could increase the relaxation errors in LiRPA.
Furthermore, in real-world applications, unmodeled dynamics such as sensor noise, latency, friction, 
and higher-order integration errors can deviate from the nominal model. 
Since nominal certificates do not inherently provide formal robustness margins against external disturbances, 
the sim-to-real gap may require including bounded parameter uncertainties and disturbance sets directly 
into the verification graph. Note that this significantly expands the verification dimensions.